\usepackage[tmargin=2cm,rmargin=1in,lmargin=1in,margin=0.85in,bmargin=2cm,footskip=.2in]{geometry}
\usepackage{amsmath,amsfonts,amsthm,amssymb,mathtools}
% \usepackage[varbb]{newpxmath}
\usepackage{xfrac}
\usepackage[makeroom]{cancel}
\usepackage{mathtools}
\usepackage{bookmark}
\usepackage{enumitem}
\usepackage{hyperref,theoremref}
\hypersetup{
	pdftitle={assignment},
	colorlinks=true, linkcolor=doc!90,
	bookmarksnumbered=true,
	bookmarksopen=true
}
\usepackage[most,many,breakable]{tcolorbox}
\usepackage{xcolor}
\usepackage{varwidth}
\usepackage{varwidth}
\usepackage{etoolbox}
%\usepackage{authblk}
\usepackage{nameref}
\usepackage{multicol,array}
\usepackage[ruled,vlined,linesnumbered]{algorithm2e}
\usepackage{comment} % enables the use of multi-line comments (\ifx \fi) 
\usepackage{import}
\usepackage{xifthen}
\usepackage{pdfpages}
\usepackage{transparent}
\usepackage{chngcntr}
\usepackage{tikz}
\usepackage{titletoc}
\usepackage{graphicx}
\usepackage{parskip}

\newcommand\mycommfont[1]{\footnotesize\ttfamily\textcolor{doc}{#1}}
\SetCommentSty{mycommfont}
\newcommand{\incfig}[2][\columnwidth]{%
	\def\svgwidth{#1}
	\import{./figures/}{#2.pdf_tex}
}

\usepackage{tikzsymbols}
\tikzset{
	symbol/.style={
			draw=none,
			every to/.append style={
					edge node={node [sloped, allow upside down, auto=false]{$#1$}}}
		}
}
\tikzstyle{c} = [circle,fill=black,scale=0.5]
\tikzstyle{b} = [draw, thick, black, -]
\tikzset{
	vertex/.style={
			circle,
			draw,
			minimum size=6mm,
			inner sep=0pt
		}
}

\renewcommand\qedsymbol{$\Laughey$}

% \usepackage{xparse}
% \usepackage{pgffor}
% \usepackage{emoji}
%
% \NewDocumentCommand{\memoji}{m o}{%
%   \IfNoValueTF{#2}
%     {\ifmmode \text{\emoji{#1}} \else \emoji{#1}\fi} % Single emoji case
%     {\foreach \index in {1,...,#2}{\ifmmode \text{\emoji{#1}} \else \emoji{#1}\fi}} % Repeated emoji case
% }

%%%%%%%%%%%%%%%%%%%%%%%%%%%%%%
% MORANDI PALETTE
%%%%%%%%%%%%%%%%%%%%%%%%%%%%%%
% Muted accents; dark headings remain legible on lightly tinted backgrounds.
\definecolor{morandiSlate}{HTML}{647681}
\definecolor{morandiSage}{HTML}{737F70}
\definecolor{morandiClay}{HTML}{A27E70}
\definecolor{morandiMauve}{HTML}{8A7888}
\definecolor{morandiSand}{HTML}{9B886D}
\definecolor{morandiRose}{HTML}{A2787B}
\definecolor{morandiInk}{HTML}{353938}

% Keep existing color names available to documents using this preamble.
\colorlet{doc}{morandiSlate}
\colorlet{myg}{morandiSage}
\colorlet{myb}{morandiSlate}
\colorlet{myr}{morandiRose}
\colorlet{myp}{morandiMauve}
\colorlet{mygr}{morandiInk}
\colorlet{myred}{morandiRose!80!black}
\colorlet{myyellow}{morandiSand}
\colorlet{notesgreen}{morandiSage}
\colorlet{mytheorembg}{morandiSlate!8!white}
\colorlet{mytheoremfr}{morandiSlate}
\colorlet{mylemmabg}{morandiClay!8!white}
\colorlet{mylemmafr}{morandiClay}
\colorlet{mypropbg}{morandiSlate!8!white}
\colorlet{mypropfr}{morandiSlate}
\colorlet{myexamplebg}{morandiSage!8!white}
\colorlet{myexamplefr}{morandiSage}
\colorlet{myexampleti}{morandiSage!75!black}
\colorlet{mydefinitbg}{morandiMauve!8!white}
\colorlet{mydefinitfr}{morandiMauve}
\colorlet{myexercisebg}{morandiSand!8!white}
\colorlet{myexercisefg}{morandiSand}

%%%%%%%%%%%%%%%%%%%%%%%%%%%%%%
% SHARED ENVIRONMENT STYLE
%%%%%%%%%%%%%%%%%%%%%%%%%%%%%%
\setlength{\parindent}{1cm}
\tcbuselibrary{theorems,skins,hooks}
\usetikzlibrary{arrows,calc,shadows.blur}
\tcbset{
  mathnote box/.style n args={2}{
    enhanced,
    breakable,
    lines before break=4,
    sharp corners,
    frame hidden,
    boxrule=0pt,
    borderline west={2pt}{0pt}{#1},
    colback=#2,
    colframe=#1,
    coltext=morandiInk,
    coltitle=#1!75!black,
    fonttitle=\bfseries\sffamily,
    fontupper=\normalfont,
    detach title,
    before upper={\ifstrempty{\tcbtitletext}{}{\tcbtitle\par\smallskip}},
    boxsep=0pt,
    left=10pt,right=10pt,top=8pt,bottom=8pt,
    before skip=10pt,after skip=10pt,
    segmentation style={solid,#1!45!white},
  },
  mathnote statement/.style n args={2}{
    mathnote box={#1}{#2},
    description font=\mdseries,
    separator sign={\enspace},
    description delimiters none,
  },
}

% Preserve environment names, numbering scopes, and label prefixes.
\newtcbtheorem[number within=section]{Theorem}{Theorem}
  {mathnote statement={mytheoremfr}{mytheorembg}}{th}
\newtcbtheorem[number within=chapter]{theorem}{Theorem}
  {mathnote statement={mytheoremfr}{mytheorembg}}{th}
\newtcbtheorem[number within=section]{Corollary}{Corollary}
  {mathnote statement={morandiMauve}{morandiMauve!8!white}}{th}
\newtcbtheorem[number within=chapter]{corollary}{Corollary}
  {mathnote statement={morandiMauve}{morandiMauve!8!white}}{th}
\newtcbtheorem[number within=section]{Lemma}{Lemma}
  {mathnote statement={mylemmafr}{mylemmabg}}{th}
\newtcbtheorem[number within=chapter]{lemma}{Lemma}
  {mathnote statement={mylemmafr}{mylemmabg}}{th}
\newtcbtheorem[number within=section]{Exercise}{Exercise}
  {mathnote statement={myexercisefg}{myexercisebg}}{th}
\newtcbtheorem[number within=chapter]{exercise}{Exercise}
  {mathnote statement={myexercisefg}{myexercisebg}}{th}
\newtcbtheorem[number within=section]{Prop}{Proposition}
  {mathnote statement={mypropfr}{mypropbg}}{th}
\newtcbtheorem[number within=chapter]{prop}{Proposition}
  {mathnote statement={mypropfr}{mypropbg}}{th}
\newtcbtheorem[number within=section]{claim}{Claim}
  {mathnote statement={morandiSage}{morandiSage!8!white}}{th}
\newtcbtheorem[number within=section]{Example}{Example}
  {mathnote statement={myexamplefr}{myexamplebg}}{ex}
\newtcbtheorem[number within=chapter]{example}{Example}
  {mathnote statement={myexamplefr}{myexamplebg}}{ex}
\newtcbtheorem[number within=section]{Definition}{Definition}
  {mathnote statement={mydefinitfr}{mydefinitbg}}{def}
\newtcbtheorem[number within=chapter]{definition}{Definition}
  {mathnote statement={mydefinitfr}{mydefinitbg}}{def}
\newtcbtheorem[number within=chapter]{wconc}{Wrong Concept}
  {mathnote statement={morandiRose}{morandiRose!8!white}}{def}

\newtcolorbox{Theoremcon}{mathnote box={mytheoremfr}{mytheorembg}}
\newcounter{questioncounter}
\counterwithin{questioncounter}{chapter}
\newtcbtheorem[use counter=questioncounter]{question}{Question}
  {mathnote statement={morandiSand}{morandiSand!8!white}}{def}
\newtcbtheorem{qstion}{Question}
  {mathnote statement={morandiSand}{morandiSand!8!white}}{def}
\newtcolorbox{solution}{mathnote box={morandiSage}{morandiSage!8!white},title=Solution}
\newtcolorbox{note}[1][]{mathnote box={morandiSand}{morandiSand!8!white},title=Note,#1}

%%%%%%%%%%%%%%%%%%%%%%%%%%%%%%
% SELF MADE COMMANDS
%%%%%%%%%%%%%%%%%%%%%%%%%%%%%%

\makeatletter
\newcommand{\thm}[2]{\begin{Theorem}{#1}{}\protected@edef\@currentlabel{\csname thetcb@cnt@Theorem\endcsname}#2\end{Theorem}}
\newcommand{\cor}[2]{\begin{Corollary}{#1}{}\protected@edef\@currentlabel{\csname thetcb@cnt@Corollary\endcsname}#2\end{Corollary}}
\newcommand{\lma}[2]{\begin{Lemma}{#1}{}\protected@edef\@currentlabel{\csname thetcb@cnt@Lemma\endcsname}#2\end{Lemma}}
\newcommand{\mer}[2]{\begin{Exercise}{#1}{}\protected@edef\@currentlabel{\csname thetcb@cnt@Exercise\endcsname}#2\end{Exercise}}
\newcommand{\mprop}[2]{\begin{Prop}{#1}{}\protected@edef\@currentlabel{\csname thetcb@cnt@Prop\endcsname}#2\end{Prop}}
\newcommand{\clm}[3]{\begin{claim}{#1}{#2}\protected@edef\@currentlabel{\csname thetcb@cnt@claim\endcsname}#3\end{claim}}
\newcommand{\wc}[2]{\begin{wconc}{#1}{}\protected@edef\@currentlabel{\csname thetcb@cnt@wconc\endcsname}\setlength{\parindent}{1cm}#2\end{wconc}}
\newcommand{\thmcon}[1]{\begin{Theoremcon}{#1}\end{Theoremcon}}
\newcommand{\ex}[2]{\begin{Example}{#1}{}\protected@edef\@currentlabel{\csname thetcb@cnt@Example\endcsname}#2\end{Example}}
\newcommand{\dfn}[2]{\begin{Definition}{#1}{}\protected@edef\@currentlabel{\csname thetcb@cnt@Definition\endcsname}#2\end{Definition}}
\newcommand{\dfnc}[2]{\begin{definition}{#1}{}\protected@edef\@currentlabel{\csname thetcb@cnt@definition\endcsname}#2\end{definition}}
\newcommand{\qs}[2]{\begin{question}{#1}{}\protected@edef\@currentlabel{\csname thetcb@cnt@question\endcsname}#2\end{question}}
\makeatother
\newcommand{\pf}[2]{\begin{myproof}[#1]#2\end{myproof}}
\newcommand{\nt}[1]{\begin{note}#1\end{note}}

\newcommand*\circled[1]{\tikz[baseline=(char.base)]{
		\node[shape=circle,draw,inner sep=1pt] (char) {#1};}}
\newcommand\getcurrentref[1]{%
	\ifnumequal{\value{#1}}{0}
	{??}
	{\the\value{#1}}%
}
\newcommand{\getCurrentSectionNumber}{\getcurrentref{section}}
\newenvironment{myproof}[1][\proofname]{%
  \begin{tcolorbox}[mathnote box={morandiSlate}{morandiSlate!8!white},title={#1}]%
  \pushQED{\qed}%
}{\popQED\end{tcolorbox}}

\newcommand{\mclm}[2]{\begin{myclaim}[#1]#2\end{myclaim}}
\providecommand{\claimname}{Claim}
\newenvironment{myclaim}[1][\claimname]{%
  \begin{tcolorbox}[mathnote box={morandiSage}{morandiSage!8!white},title={#1}]%
}{\end{tcolorbox}}
\newenvironment{iclaim}[1][\claimname]{\begin{myclaim}[#1]}{\end{myclaim}}
\newcommand{\iclm}[2]{\begin{iclaim}[#1]#2\end{iclaim}}

\newcounter{mylabelcounter}

\makeatletter
\newcommand{\setword}[2]{%
	\phantomsection
	#1\def\@currentlabel{\unexpanded{#1}}\label{#2}%
}
\makeatother

% deliminators
\DeclarePairedDelimiter{\abs}{\lvert}{\rvert}
\DeclarePairedDelimiter{\norm}{\lVert}{\rVert}

\DeclarePairedDelimiter{\ceil}{\lceil}{\rceil}
\DeclarePairedDelimiter{\floor}{\lfloor}{\rfloor}
\DeclarePairedDelimiter{\round}{\lfloor}{\rceil}

\newsavebox\diffdbox
\newcommand{\slantedromand}{{\mathpalette\makesl{d}}}
\newcommand{\makesl}[2]{%
	\begingroup
	\sbox{\diffdbox}{$\mathsurround=0pt#1\mathrm{#2}$}%
	\pdfsave
	\pdfsetmatrix{1 0 0.2 1}%
	\rlap{\usebox{\diffdbox}}%
	\pdfrestore
	\hskip\wd\diffdbox
	\endgroup
}
\newcommand{\dd}[1][]{\ensuremath{\mathop{}\!\ifstrempty{#1}{%
			\slantedromand\@ifnextchar^{\hspace{0.2ex}}{\hspace{0.1ex}}}%
		{\slantedromand\hspace{0.2ex}^{#1}}}}
\ProvideDocumentCommand\dv{o m g}{%
	\ensuremath{%
		\IfValueTF{#3}{%
			\IfNoValueTF{#1}{%
				\frac{\dd #2}{\dd #3}%
			}{%
				\frac{\dd^{#1} #2}{\dd #3^{#1}}%
			}%
		}{%
			\IfNoValueTF{#1}{%
				\frac{\dd}{\dd #2}%
			}{%
				\frac{\dd^{#1}}{\dd #2^{#1}}%
			}%
		}%
	}%
}
\providecommand*{\pdv}[3][]{\frac{\partial^{#1}#2}{\partial#3^{#1}}}
%  - others
\DeclareMathOperator{\Lap}{\mathcal{L}}
\DeclareMathOperator{\Var}{Var} % varience
\DeclareMathOperator{\Cov}{Cov} % covarience
\DeclareMathOperator{\E}{E} % expected

% Since the amsthm package isn't loaded

% I prefer the slanted \leq
\let\oldleq\leq % save them in case they're every wanted
\let\oldgeq\geq
\renewcommand{\leq}{\leqslant}
\renewcommand{\geq}{\geqslant}

%%%%%%%%%%%%%%%%%%%%%%%%%%%%%%%%%%%%%%%%%%%
% TABLE OF CONTENTS
%%%%%%%%%%%%%%%%%%%%%%%%%%%%%%%%%%%%%%%%%%%

\contentsmargin{0cm}
\titlecontents{chapter}[3.7pc]
{\addvspace{30pt}%
	\begin{tikzpicture}[remember picture, overlay]%
		\draw[fill=doc!60,draw=doc!60] (-7,-.1) rectangle (-0.9,.5);%
		\pgftext[left,x=-3.7cm,y=0.2cm]{\color{white}\Large\sc\bfseries Chapter\ \thecontentslabel};%
	\end{tikzpicture}\color{doc!60}\large\sc\bfseries}%
{}
{}
{\;\titlerule\;\large\sc\bfseries Page \thecontentspage
	\begin{tikzpicture}[remember picture, overlay]
		\draw[fill=doc!60,draw=doc!60] (2pt,0) rectangle (4,0.1pt);
	\end{tikzpicture}}%
\titlecontents{section}[3.7pc]
{\addvspace{2pt}}
{\contentslabel[\thecontentslabel]{2pc}}
{}
{\hfill\small \thecontentspage}
[]
\titlecontents*{subsection}[3.7pc]
{\addvspace{-1pt}\small}
{}
{}
{\ --- \small\thecontentspage}
[ \textbullet\ ][]

\makeatletter
\renewcommand{\tableofcontents}{%
	\chapter*{%
	  \vspace*{-20\p@}%
	  \begin{tikzpicture}[remember picture, overlay]%
		  \pgftext[right,x=15cm,y=0.2cm]{\color{doc!60}\Huge\sc\bfseries \contentsname};%
		  \draw[fill=doc!60,draw=doc!60] (13,-.75) rectangle (20,1);%
		  \clip (13,-.75) rectangle (20,1);
		  \pgftext[right,x=15cm,y=0.2cm]{\color{white}\Huge\sc\bfseries \contentsname};%
	  \end{tikzpicture}}%
	\@starttoc{toc}}
\makeatother

% Append by Friedforks
\newcommand{\inlinecaption}[1]{\par\smallskip\centering\textbf{#1}\par\smallskip}

\newcommand{\centeredimage}[2][]{%
	\begin{center}
		\includegraphics[#1]{#2}%
	\end{center}%
}

\NewDocumentCommand{\code}{m m m}{
  \begin{tcolorbox}[mathnote box={morandiSlate}{morandiSlate!8!white},title={#1}]
    \begin{minted}{#2}
      #3
    \end{minted}
  \end{tcolorbox}
}
